\subsection{KV Master: Request-Specific Backup Decisions}
\label{sec:kv-placement}
\label{sec:kv-master}

KV Master (3) selects recovery tiers, each request's Peer depth $L_r$,
and layer-local hot sets. Recovery Manager fixes the retry/Peer target;
Adapter publishes its window-level capacity $B_p$. Requests may use
different depths within the same capacity window.

\paragraph{Decision events.}
The Master establishes policy at arrival, offers sealed state after prefill,
and reviews a request every additional $Z$ generated tokens. $Z$ is a
block-aligned system configuration, not an Adapter output. Significant
slack, budget, context, or copy-state changes can trigger an earlier review.
A decision does not imply a completed checkpoint: the replay tail extends
beyond the latest committed frontier and can exceed $Z$ when copying lags.
Copy queues have configured byte and operation limits. Coalesce superseded
unstarted plans, retain dependencies of in-flight or committed versions,
and skip a new backup attempt if it cannot be reserved. Report the resulting
commit lag; the trigger interval supplies no durability bound.
Sealing only updates eligibility. Hot-block statistics use the preceding
512 decode steps, independently of checkpoint and capacity windows.

\paragraph{Choose protection tiers.}
Let $T_{\mathrm{rc}}$ denote prefix recomputation, $T_{\mathrm{rp}}$ known-tail
replay, $W_0$ exposed DRAM-loading wait, and $R_{\mathrm{floor}}$ detection,
rerouting, and readmission time. Joint context/tail/load profiles give
\begin{equation}
\begin{aligned}
 R_{\mathrm{base}}&=R_{\mathrm{floor}}+T_{\mathrm{rp}},\\
 R_R&=R_{\mathrm{base}}+T_{\mathrm{rc}},\qquad
 R_D=R_{\mathrm{base}}+W_0.
\end{aligned}
\label{eq:recovery-base}
\end{equation}
$R_R$ and $R_D$ use the weighted conditional means of the same
future-state samples defined below; the tier test is a planning heuristic,
not a deadline guarantee. Tier R needs no backup if $\sigma_r\ge R_R$.
Otherwise choose D when $R_D<R_R$ and $\sigma_r\ge R_D$; remaining
requests compete for Peer as tier P. There is no assumption that loading
is faster than recomputation. The non-Peer fallback chooses the cheaper
feasible route, preferring R on equal costs; a denied Peer increment does
not force a slower DRAM restore. DRAM protection requires its own admission
and commit (\S\ref{sec:kv-runtime}). Missing profiles request conservative
DRAM protection when feasible, without claiming it is faster; failed
admission uses an older valid checkpoint or R. Slack orders protection,
not retry execution. At recovery, refresh costs and availability; overlapping
cost intervals retain the feasible current route unless dominance is clear.

\paragraph{K counts resident blocks, not selector outputs.}
A protection object is an immutable sealed \emph{(layer, block)} including
all KV heads; the trace study uses 64-token blocks. Let $H_{r,l}$ be
request $r$'s published, committed Peer set at layer $l$. We define
\begin{equation}
 K_{r,l}=|H_{r,l}|.
 \label{eq:hot-block-count}
\end{equation}
Thus $K$ counts Peer-resident blocks of one request and layer. It is not
the model's per-query top-$k$, a hit-rate guarantee, or a count of copies
still in flight. Physical allocation also charges retiring buffers,
reservations, and selector state; shared copies are counted once per pool.

A separately configured quota determines the desired hot-set size.
The configuration declares either a fixed block limit or
$\lceil h_l n_r\rceil$ blocks for $n_r$ sealed historical blocks.
Mandatory sealed sink/window blocks are included in that limit: if they
alone exceed it, request a larger explicitly charged quota or decline
that layer's Peer admission, never silently drop native requirements.
The rule is
fixed in each comparison; this paper does not claim an optimal quota.
Remaining slots rank native selection counts, then recency and stable
block ID. Optional blocks leave after two consecutive statistics windows
outside the target set. Required auxiliary state is budgeted separately.
Before 512 steps are available, use the observed prefix; with no observations,
rank mandatory blocks first and then recency and stable ID. This cold-start
rule has no demonstrated prediction accuracy.
For a replay of $N$ tokens, demand is the deduplicated union of native
selections across queries and KV groups, not $N$ copies of a single top-$k$.

\paragraph{Preview replacement as well as admission.}
A candidate depth $d$ offers target hot sets in layers $1{:}d$;
this prefix family restricts the search, rather than proving that arbitrary
layer placements cannot do better.
Its plan retains intersections, demotes owned blocks leaving the set,
and admits new blocks. It cannot evict another request's protected copy.
A shared object loses only the departing owner's reference.
Dirty demotions, new transfers, and temporary space all enter the preview;
slots become reusable only after actual release
(\S\ref{sec:kv-runtime}). Shrinking $L_r$ can therefore retire the owner's
deeper copies. An unchanged depth can still replace block identities.
Uniform, fully admitted layers contain $L_r K$ layer-block objects;
partial admission is reported through the actual sets $H_{r,l}$.

\paragraph{A reproducible depth objective.}
All candidates start from the same target snapshot $S$ and use the same
planning horizon $H$, future-state samples $j$, and normalized weights $w_j$,
conditional on a source failure within $H$. Current slack is held fixed
within this local comparison and refreshed at the next policy event.
A sample specifies committed history $M_j$, replay tail $N_j$, available
copy time, and foreground/retry load. Samples come from a declared,
versioned profile, including commit lag; they do not assume $N_j=Z$ or
a uniform checkpoint phase. Co-located requests share the same source-event
samples and contended target resources. For a configured Poisson model,
$\pi_r=1-\exp(-\lambda_r H)$. Other risk models must provide the corresponding
probability over the same $H$; a failure rate is not itself a probability.

Let $G_d$ be the admitted-and-retired plan from the preview.
$A_r(d)$ is the predicted aggregate change in foreground request-delay
seconds, relative to keeping the current placement at $S$. The cohort
contains active requests and arrivals within $H$, followed to completion
under matched arrivals and a held policy, including completions after $H$.
It includes replication, writeback, transient occupancy, and reduced or
increased admission pressure, and may be negative. Let $R_r(d,j)$ be
failure-to-ready latency for sample $j$, crediting only replicas predicted
to commit before its failure. The Master minimizes
\begin{equation}
 J_r(d)=A_r(d)+\pi_r\sum_j w_j
          \max\{0,R_r(d,j)-\sigma_r\}.
 \label{eq:depth-score}
\end{equation}
Both terms have units of delay seconds. This is a local protection
heuristic in slack order; it neither sums to workflow makespan nor
optimizes all later requests' allocations.
For example, two equally critical branches on one failed source each
need 10\,s to recover. Protecting only one reduces their summed local
penalties from 20 to 10\,s but leaves workflow delay at 10\,s. Source-event
DAG replay and a feasible small-DAG oracle are required to assess this gap.

\paragraph{Profile interface and validity.}
A profile query takes the model/selector and weight revisions, KV layout,
block size, hardware/link topology, $S$, $G_d$, the shared state samples,
and the configured copy scheduler. It returns $A_r(d)$ and $R_r(d,j)$
with uncertainty intervals, a generation ID, support ranges, and expiry.
$A$ is obtained from matched normal-service traces with identical arrivals;
$R$ uses matched recovery traces or a declared, calibrated event model.
Transfer and compute overlap are evaluated on one timeline; adding their
full durations would double-count hidden transfer. Rows state whether
they are measured or modeled. Missing measurements are not filled with
the Figure~\ref{fig:design-pipeline} example.

The finite depth grid includes zero and the current depth, and an explicit
\emph{keep} plan preserves current block identities. The lowest score wins;
ties prefer smaller depth, then stable plan ID. Propagating the profile's
intervals through (\ref{eq:depth-score}) gives lower and upper score bounds.
A change is accepted only if its upper bound plus a configured margin
is below the keep plan's lower bound. Unsupported candidates are excluded.
An invalid generation or missing, expired, or drift-flagged keep profile
freezes discretionary expansion and replacement; new requests attempt
DRAM admission, with the stated fallback if it fails.
Existing protection is retained only while feasible:
Adapter contraction and correctness-driven retirement still take precedence.
A generation is invalidated after a configured number of consecutive
relative-error violations on comparable observed operations; unmatched
observations do not count as validation. Thresholds, grid, horizon,
risk model, quota rule, sample weights, and calibration version are
recorded configuration, not inferred experimental constants.
At zero failure probability, only $A$ affects the score; costly expansion
has no recovery justification, while a zero-cost keep plan need not move.
Uncertain route crossings do not justify an unprofiled migration.

\begin{figure}[t]
  \raggedright
  \refstepcounter{designalgorithm}\label{alg:placement}
  \noindent\textbf{Algorithm \thedesignalgorithm: Per-request backup policy}
  \par\smallskip
  \noindent Input: requests, slack, fixed $Z$, profiles, targets, budgets.
  Output: admitted tiers, Peer plans, and committed $H_{r,l}$.
  \par\smallskip
  \begin{tabularx}{\columnwidth}{@{}r@{\quad}X@{}}
    \toprule
    1 & Snapshot copies, leases, budgets, and reservations as $S$ \\
    2 & \textbf{for} affected $r$ in $(\sigma_r,\mathrm{ID}_r)$ order: \\
    3 & \quad Choose R/D/P and the feasible non-Peer fallback \\
    4 & \quad For D/P, request DRAM; report coverage/failure \\
    5 & \quad Rank layer-local candidates under the quota rule \\
    6 & \quad Form keep and depth plans; R/D target depth is zero \\
    7 & \quad Preview each plan independently from the same $S$ \\
    8 & \quad Query valid profiles; score by (\ref{eq:depth-score}) \\
    9 & \quad Choose a stable improvement, or feasible fallback \\
    10 & \quad Queue its retains, demotions, and new admissions \\
    11 & \quad Update $S$ with only this chosen plan \\
    12 & Atomically recheck reservations before physical execution \\
    13 & Publish only committed residency and actual released bytes \\
    \bottomrule
  \end{tabularx}
  \Description{The Master chooses per-request layers and hot sets using
    a shared-sample cost objective. Only chosen plans change the virtual
    allocation snapshot; execution rechecks space and safe retirement.}
\end{figure}

The preview is not a reservation. Actual grants may be smaller. The
directory reports committed $K_{r,l}$ and frontiers; a mismatch triggers
another review. Sorting $n$ requests takes $O(n\log n)$ time. Candidate
work scales with the grid, sample count, and affected blocks. Profile
calibration and accuracy remain evaluation work.

\paragraph{Keep selection and placement separate.}
Recovery restores model-specific selector state before invoking the
unchanged native selector. The selector may name any causally eligible
historical or recomputed KV, including DRAM-only objects. The Connector
resolves logical IDs to committed replicas, fetches actual misses, and
returns readiness events; attention waits for every required input.
Forecast mistakes add traffic or waiting, not a new attention approximation.
Tier R and missing committed prefixes use dependency-correct recomputation
of logged tokens. A batched replay is eligible only if its native selector
and state evolution are equivalent to the original token-wise path under
the declared numerical contract; otherwise replay token by token. A native
prefill gate or an unresolved numerical diagnostic does not establish this
equivalence. Compatible weights, layout, and operators are required;
future stochastic samples need not be bitwise identical across floating-point
execution schedules.

\begin{figure}[t]
  \centering
  \includegraphics[width=\columnwidth]{figures/design-single/design-layer-pipeline.pdf}
  \caption{Modeled recovery pipelines: MiniCPM4-8B, 32 layers, 32K
    history, 64-token known tail; $C=2.87$\,ms/layer (A100-derived average)
    and $D=5.37$\,ms at an assumed 6.25\,GB/s. In (d), 147\,MiB of Peer
    residency removes preload wait. Totals end at tail-KV reconstruction.
    The homogeneous, preemptible-copy model excludes detection,
    queuing, and interference.}
  \label{fig:design-pipeline}
  \Description{Four timelines compare token-block transfer, full-layer
    pipelining, hot preloads with demand fetch, and Peer-resident hot sets.}
\end{figure}

Figure~\ref{fig:design-pipeline} illustrates one fixed-tail example:
layer pipelining, hot preloads, and seven resident hot sets reduce exposed
wait from 172 to 83, 44, and 20\,ms. It does not calibrate
(\ref{eq:depth-score}), establish multi-request WCT gains, or make seven
layers appropriate for every checkpoint state.
Here $C$ is a five-repeat median of a batch-one, 64-token native-path
forward over resident history on an A100 80GB PCIe, divided by 32 layers;
individual layer timings were not measured. The assumed bandwidth is a
25\,GB/s pool link equally shared by four recoveries, not measured H2D.
With 1\,KiB/token/layer, the Peer payload is
$7\times0.65845\times32$\,MiB $\simeq147.49$\,MiB, excluding auxiliary
state and temporary buffers. Fractions are trace-derived means, not an
integer per-layer grant. The model leaves 38.69\,ms of cold-copy work
after tail readiness; subsequent token latency requires separate evaluation.
